\section{Step 41 Locate DC without centering INTRA samples}\label{h263-transform-s01}
Begin with the flat forward observation. Predict the DC value from the basis
normalization before writing a loop: each dimension contributes its constant
basis factor, giving DC eight times the constant sample. A block filled with
128 therefore has DC 1024. Its AC terms vanish because their signed basis values
cancel over a constant row or column.

Keep this distinct from an INTER residual that happens to be zero. The transform
entry accepts signed sample values, and its caller decides whether they represent
INTRA samples or prediction residuals. Do not insert an unrequested level shift
inside the transform to make the two cases look alike.

\paragraph{Observe} Flat forward output begins with 1024. Zero input has all
sixty-four outputs exactly zero in either direction.
\Needspace{8\baselineskip}
\paragraph{Check} What should happen to a constant INTRA block of value 128 here?
\begin{itemize}
\item[A.] Subtract 128 and force all coefficients to zero.
\item[B.] Produce DC 1024 with zero AC coefficients.
\item[C.] Treat each input byte as an already coded coefficient.
\end{itemize}
